\section{误差如何产生与放大}
\label{sec:08-Numerical-Analysis/01-Floating-Point-and-Error/02}

有一段代码，几乎每个人都写过：算一组数的方差，用课本上那条公式

\[
\Var(X)=\E[X^2]-(\E X)^2 .
\]

拿它去算一批传感器读数，比如围绕 \(1{,}000{,}000\) 上下浮动 \(0.01\) 的一万个数。跑出来的方差是负的。

方差是平方的平均，不可能为负。可这段代码没有任何逻辑错误，公式也没抄错。它只是把两个都在 \(10^{12}\) 量级的数相减，而正确答案在 \(10^{-4}\) 量级——中间隔了十六个数量级，正好是 binary64 的全部家当。

\subsection{相减不制造误差，它只是揭穿误差}

回到上一节那张数轴图。两个数值接近的浮点数，它们的高位有效数字是一样的，差异只藏在末尾几位。把它们相减，高位整齐地抵消掉，剩下的那几位低位被重新规格化、左移到最高位——于是原本躲在有效位末端、只占相对误差 \(u\) 的那点扰动，一跃成为结果的主导成分。

\begin{center}
\begin{tikzpicture}[x=0.5cm,y=1cm,font=\small]
  \node[left,font=\footnotesize] at (-0.25,2.22) {\(a\)};
  \node[left,font=\footnotesize] at (-0.25,1.32) {\(b\)};
  \node[left,font=\footnotesize] at (-0.25,0.32) {\(a-b\)};
  \foreach \i in {0,...,9}{
    \draw[fill=black!16] (\i,2) rectangle (\i+1,2.45);
    \draw[fill=black!16] (\i,1.1) rectangle (\i+1,1.55);
  }
  \foreach \i in {10,...,15}{
    \draw[fill=blue!22] (\i,2) rectangle (\i+1,2.45);
    \draw[fill=orange!30] (\i,1.1) rectangle (\i+1,1.55);
  }
  \foreach \i in {0,...,5}
    \draw[fill=green!20] (\i,0.1) rectangle (\i+1,0.55);
  \foreach \i in {6,...,15}
    \draw[fill=red!14] (\i,0.1) rectangle (\i+1,0.55);
  \node[font=\footnotesize] at (11,0.32) {填进来的全是无信息位};
  \draw[black!55] (0,2.75) -- (0,2.9) -- (10,2.9) -- (10,2.75);
  \node[above,font=\footnotesize,black!75] at (5,2.9) {两数完全相同的高位};
  \draw[black!55] (0,-0.15) -- (0,-0.3) -- (6,-0.3) -- (6,-0.15);
  \node[below,font=\footnotesize,black!75] at (3,-0.3) {真正携带信息的只剩这几位};
\end{tikzpicture}

\emph{每个方格代表一个有效位。相减后高位成片归零，结果左移规格化，尾部只能由零或舍入噪声补齐。位数看上去没少，可信位少了十位。}
\end{center}

这张图解释了为什么这类事故被叫作灾难性抵消（catastrophic cancellation），也解释了为什么``提高精度''往往救不了它。信息不是在减法那一步丢的，是在两个大数被存进来的时候就已经不在了；减法只是把这件事暴露出来。用 binary64 重跑一遍上面的方差代码，负号可能消失，但有效位照样只剩个位数——把量级换成 \(10^{16}\)，负方差又会回来。

\subsection{同一个数学式，几种写法，几种命运}

既然损失发生在表达式的某一步，换一条计算路径就有救。

小 \(x\) 时计算 \((\sqrt{1+x}-1)/x\) 是最干净的例子。数学上它趋于 \(1/2\)，可 \(\sqrt{1+x}\) 与 1 在前十几位完全一样，直接算等于让上面那张图重演一遍。有理化之后

\[
\frac{\sqrt{1+x}-1}{x}=\frac{1}{\sqrt{1+x}+1},
\]

右式的分母是两个正数相加，没有任何抵消，取 \(x=10^{-16}\) 也能稳稳给出 \(0.5\)。同一条思路对付 \(\sqrt{x^2+1}-x\)（大 \(x\)）和 \(1-\cos x\)（小 \(x\)，改写成 \(2\sin^2(x/2)\)）一样有效。

二次方程求根公式则示范了另一种手法。当 \(b^2\gg 4ac\) 时，\(\sqrt{b^2-4ac}\approx\abs{b}\)，于是

\[
x_{1,2}=\frac{-b\pm\sqrt{b^2-4ac}}{2a}
\]

两个符号里必有一个是灾难：它让 \(-b\) 和一个几乎等于 \(b\) 的数相减。绝对值较大的那个根没问题，绝对值较小的那个根会被抵消吃掉大半有效位。补救办法不是有理化，而是换分支——先算安全的那个根

\[
q=-\tfrac12\left(b+\sign(b)\sqrt{b^2-4ac}\right),
\]

再用根与系数的关系 \(x_1x_2=c/a\) 把另一个根接出来：

\[
x_1=\frac{q}{a},\qquad x_2=\frac{c}{q}.
\]

实现时要把 \(\sign(0)\) 显式定成 \(1\)，否则 \(b=0\) 会让 \(q\) 变成零。这里没有改动任何数学，只是不让任何一个根经由相近数相减产生。

方差的补救属于同一类。把``先各自求和、最后相减''换成边扫描边更新的 Welford 递推：维护均值 \(\bar x_n\) 与偏差平方和 \(M_n\)，每来一个新样本做

\[
\begin{aligned}
\bar x_n&=\bar x_{n-1}+\frac{x_n-\bar x_{n-1}}{n},\\
M_n&=M_{n-1}+(x_n-\bar x_{n-1})(x_n-\bar x_n).
\end{aligned}
\]

式子里出现的减法都是``样本减当前均值''，两者本来就在同一量级，差值不小，没有抵消可言；\(M_n\) 也保证非负，负方差从结构上就不可能出现。批归一化、在线统计、流式监控里的方差实现基本都走这条路，理由就在这里。

\subsection{小量在长求和里被一个个吞掉}

抵消是大数相减的问题，求和则有一个镜像的问题：小数加不进大数。

设部分和已经累到 \(S\)，再加一个 \(x\)。若 \(\abs{x}\) 小于 \(S\) 在该处刻度的一半，\(S+x\) 舍入之后还是 \(S\)，这一项就等于没加。一百万个 \(10^{-8}\) 顺序相加，理论和是 \(10^{-2}\)；实际算下来，部分和一旦长到某个尺度，后面的项就开始逐个失效，最终结果可能差出一个数量级，而不是几个 ulp。

对策有三层，代价递增。按绝对值从小到大排序再加，让小量先抱团长大；成对求和把累加组织成二叉树，每次合并规模相近的两块，误差随规模的增长从 \(O(n)\) 降到 \(O(\log n)\)，NumPy 的 \texttt{sum} 就用它；Kahan 补偿求和额外维护一个变量，专门记录每次加法被舍掉的零头，下一次补回去。但这三招都有同一个前提：最终和本身不是一堆正负大数抵消后的微小残差。如果是，那就回到了上一小节的处境，求和技术无能为力。

\subsection{三种误差，量级常常差得很远}

到这里出现的都是舍入误差。它不是唯一的来源，甚至常常不是最大的那个。

建模误差来自问题的提法：假设传感器噪声是高斯的，假设特征与标签是线性关系，假设摩擦力可以忽略。哪怕后面每一步都精确无误，答案也只对那个被简化过的世界成立。

截断误差来自用有限过程替代无限过程：\(e^x\) 只取到 Taylor 展开的前三项，导数用差商 \((f(x+h)-f(x))/h\) 近似，微分方程用有限步 Euler 法推进。哪怕有一台位数无限的计算机，这类误差照样存在。

舍入误差才是本单元前两节的主角，量级由 \(u\) 决定。

\begin{center}
\begin{tikzpicture}[x=0.52cm,y=1cm,font=\small]
  \draw[->,black!55] (0,-0.35) -- (15.6,-0.35);
  \foreach \xa/\lab in {0/{\(10^{-16}\)},4/{\(10^{-12}\)},8/{\(10^{-8}\)},12/{\(10^{-4}\)},15/{\(10^{-1}\)}}{
    \draw[black!45] (\xa,-0.45) -- (\xa,-0.25);
    \node[below,font=\footnotesize,black!70] at (\xa,-0.45) {\lab};
  }
  \draw[fill=red!25,draw=black!40] (0,1.85) rectangle (14,2.3);
  \draw[fill=orange!30,draw=black!40] (0,1.05) rectangle (10,1.5);
  \draw[fill=blue!25,draw=black!40] (0,0.25) rectangle (0.35,0.7);
  \node[right,font=\footnotesize] at (14.2,2.08) {建模};
  \node[right,font=\footnotesize] at (10.2,1.28) {截断};
  \node[right,font=\footnotesize] at (0.6,0.48) {舍入};
  \node[font=\footnotesize,black!70] at (7.8,-1.05) {误差大小};
\end{tikzpicture}

\emph{一个典型的科学计算任务里三种误差的相对量级（横轴为对数刻度）。把舍入误差从 \(10^{-16}\) 压到 \(10^{-19}\) 需要换用四倍精度、付出十倍时间，而总误差纹丝不动——它由最上面那根条决定。}
\end{center}

图里那根最短的条常常是被优化得最起劲的一根。用双精度重跑一个建模误差在 \(1\%\) 的仿真，或者为一个 Euler 法的一阶截断误差去纠结加法顺序，收益都接近于零。反过来，截断误差与舍入误差还会互相拉扯：数值微分的步长 \(h\) 越小，截断误差越小，但 \(f(x+h)-f(x)\) 的抵消越严重，总误差先降后升，最优步长落在两条曲线的交点上。这件事在\hyperref[ch:064]{``数值微分与积分：同一份采样，为何难度相反''}里会被完整算一遍。

\subsection{放大倍数写在导数里}

前面所有例子都能收进同一个框架。输入带着扰动 \(\Delta x\) 进入函数 \(f\)，一阶展开给出

\[
f(x+\Delta x)-f(x)\approx f'(x)\,\Delta x,
\]

多元情形下换成 Jacobian，\(\Delta y\approx J_f(x)\Delta x\)。导数就是那台扩音器：它大，输入末位的抖动被放成输出的显著偏差；它小，噪声被压下去。Jacobian 还多做一件事——不同方向的扰动被放大的倍数可以差出许多个数量级，某些方向几乎被滤干净，另一些方向被剧烈拉伸。

于是计算误差不是 \(u\) 乘一个固定常数，它取决于你站在函数的哪一点上。而这个放大倍数完全由 \(f\) 和 \(x\) 决定，跟你用什么算法、什么语言、什么精度都没关系。下一节给它一个名字，并说明为什么它划出了任何算法都无法逾越的下限。
